\section{A context-dependent contract for the existing increment of N}
\label{sec:state18-context}

This addendum extends the generic increment argument without modifying the
frozen runtime evidence packet. All state numbers below are old Routing245
labels. Let the dispatcher pointer $p$ be nonempty, contain no seven
consecutive zeros, and have at least three blank gap cells before the header.
Registers are arbitrary natural numbers, and the spare suffix is an arbitrary
nonempty finite list. The external decoder is transparent: it preserves the
pointer and returns to the same atomic entry on internal returns with an
unchanged pointer.

\paragraph{Theorem.}
The existing state 18 correctly increments logical register $N$, appends a
spare of value one, and performs the prescribed pointer successor whenever
\[
p\in\{0,1\}^{*}1\;\cup\;\{0,1\}^{*}10\;\cup\;\{0\}.
\]
The theorem uses the same arbitrary-register mark and increment induction
contracts as the six-entry increment-orientation lemma. It adds no machine
state and changes no transition.

\paragraph{Literal zero-pulse reduction.}
Put the first gap cell at coordinate zero; the last pointer bit is then at
$-1$. On a zero gap cell, the relevant literal rows are
\[
\begin{array}{lll}
18/0=0L27,&27/1=1R39,&39/0=0R57,\\
27/0=0L32,&32/1=1R46,&46/0=0R39.
\end{array}
\]
If $p$ ends in one, the first line gives the exact trace
\[
(18,0)\longrightarrow(27,-1)\longrightarrow(39,0)
       \longrightarrow(57,1).
\]
It preserves every tape cell. If $p$ ends in $10$, the exact trace is
\[
\begin{split}
(18,0)&\longrightarrow(27,-1)\longrightarrow(32,-2)\\
&\longrightarrow(46,-1)\longrightarrow(39,0)
 \longrightarrow(57,1).
\end{split}
\]
Again every tape cell is preserved. The singleton pointer $0$ follows the
second trace, reading the marker one at coordinate $-2$. In each case the
result is exactly that of the generic zero pulse $0R57$, with two or four
extra literal steps. No cell outside the displayed window is read or
written during these replacements.

\paragraph{Induction through the complete increment.}
The one row is already
\[
18/1=0R28,
\]
and selector 28 is the inherited mark entry for physical register six,
which represents logical register $N$. The generic proof requires only
a zero pulse ending in state 57 one cell to the right, and a one pulse
clearing that gap cell and entering the selected mark routine. The local
traces establish precisely those premises for state 18.

The first zero pulse fills the gap through state 42 and returns to the
unchanged pointer. The one pulse marks physical register six and returns
again to that unchanged pointer. Therefore the suffix hypothesis remains
true at the second zero pulse. That pulse clears the second temporary gap
cell through $57/1=0R73$ and invokes the generic marked increment. The
inherited register induction increases exactly register six, appends
$B(1)$, and returns through states 48 and 63. State 7 then performs the
binary carry, and states 10, 14, 20, and 2 restore the canonical dispatcher
frame, including the overflow marker displacement. The marker scanner
preconditions are exactly those in the generic proof: the only scan before
the increment has at most one extra zero next to the original pointer.

Both uses of the substituted zero pulse see the original unchanged pointer,
so no new suffix induction hypothesis is needed after binary carry. All
remaining phases are unchanged, and their termination follows from the same
unary and finite-list inductions. Relative to the hypothetical generic
entry $0R57,0R28$, this execution adds exactly four steps when $p$ ends in one
and eight when it ends in $10$ or is the singleton $0$. This proves the
theorem.

\subsection{Derived incoming masks for grammar words}

Write $C$ for the fine token implemented by state 18. A transparent binary
rule $w\to ab$ enters its left child after reading a zero and its right
child after reading a one. Let $M(w)\subseteq\{0,1\}$ be the incoming
last-bit values for which the word's interface is admitted. The argument
below is conditional on the declared interfaces of all other atomic
entries; observed source masks alone are not proofs of those interfaces.

Keep the ordinary mask $M(C)=\{1\}$, but use the stronger theorem when $C$
is the immediate left child. If the containing word was entered after a
one, that call to $C$ has pointer suffix $10$. Hence the recurrence is
\[
M(w)=
\begin{cases}
\varnothing,&1\notin M(b),\\
\{1\},&1\in M(b)\text{ and }a=C,\\
\{0,1\},&1\in M(b),\ a\ne C,\ 0\in M(a),\\
\varnothing,&\text{otherwise}.
\end{cases}
\]
The right-child check already accepts $b=C$, since its ordinary mask
contains one. A nonterminal left child must itself admit incoming zero;
a nonterminal right child must admit incoming one. In particular, a
right-only derived word cannot be embedded on a left edge merely because
its own parent is entered from a one: the child immediately sees the left
edge's zero. This distinction is essential.

For an acyclic grammar, induction on the word dependency graph proves
the recurrence sound relative to its atomic premises. A direct $C$ on the
left introduces the sole new dependence on the containing word's incoming
bit. All other child interfaces are checked on the actual edge bit. A word
whose mask is empty is rejected. Roots with unknown entry bit, including
conservatively treated routing roots, require mask $\{0,1\}$. Requiring
both bits even at the absolute root deliberately forgoes the exceptional
singleton-marker case; this is conservative and cannot introduce an
unsound grammar.

\paragraph{Implementation review.}
The independently inspected implementation in
\nolinkurl{anneal_fine_contextual18.cpp} represents masks by bits one and two
for incoming zero and one. Its \texttt{derived\_mask} and
\texttt{context\_ok} functions implement the displayed recurrence, require
both bits at every declared root and routing root, and reject empty masks
at all retained grammar words. These checks run on the seed and after
closure before an accepted mutation. Proposal routines may provisionally
permit $C$ on a left edge, because the final derived-mask check enforces
the contextual condition.

\paragraph{Verification boundary.}
The exact local zero-pulse trace has an executable row checker in this
addendum. The unbounded increment argument uses the inherited induction
contracts. The grammar recurrence is a compositional rule relative to
atomic contracts, not a proof of a whole new machine or of its blank
initializer. Every resulting table still requires its actual whole-frame
and blank-start verification. No new proof-kernel endpoint is claimed.
